\section{Limitations and Ethics}

\subsection{Coverage limits}

The corpus should not be described as a complete census of all Slovene Bluesky posts ever published. It covers currently recoverable public posts from discovered Slovene-linked accounts; deleted posts are excluded by design. The false-negative analysis (Section~\ref{sec:validation}) estimates an 8.5\% miss rate within the pool of excluded untagged posts from known authors that met the project's minimum-text rule; the true miss rate across the full user population is unknown and likely higher.

The methodology uses post-level rather than thread-level language classification. Slovene-linked users may post in English or other languages within otherwise Slovene conversations, and such replies are not automatically treated as Slovene corpus material.

\subsection{Open data and reproducibility}

One of the striking features of building this corpus is how open the data is. The AT Protocol was designed from the outset as a public infrastructure: there are no API keys required to access the public firehose, no rate limits on repository access, and no centralised gatekeeper deciding who can collect what. Anyone with a working internet connection and modest programming skills can reproduce the collection described here. For a small-language community like Slovene, this is an unusual and genuinely valuable property: it makes corpus construction feasible without institutional data-access agreements, and it allows the methodology to be fully transparent and reproducible.

This openness is what makes the research possible, and it is worth recognising as a design choice that distinguishes ATProto from its predecessors.

\subsection{Privacy and identifiability}

The same openness that makes collection easy raises a concern worth stating plainly. The Slovene Bluesky community, at the time of collection, consists of 432 identified authors. This is a small number. Someone reading this paper, or the corpus, could in most cases identify the individuals behind the accounts with little effort --- handles are publicly visible, many users link their real names or professional identities through custom domains, and the posting history of any individual is freely accessible through DID resolution. The corpus therefore represents not just a dataset but a detailed longitudinal record of specific, identifiable people.

This is not unique to Bluesky or to this corpus: it is a structural feature of any public social-media corpus at this scale. It is worth noting as a reminder that \textit{public} does not mean \textit{anonymous}. Users who post on Bluesky under their real names or with professionally linked handles are making a choice to be visible, but they may not have anticipated that their posting history would be compiled into a citable research corpus. The study follows standard practice by not reproducing full post texts in this article and recommending that any redistribution of the corpus take a conservative form, such as providing stable identifiers rather than full-text dumps, to limit the ease of re-identification for users who may later wish to reduce their public footprint.

Only public content was collected. Even so, reproducibility is better supported through code, decision rules, validation documentation, and aggregate statistics than through raw full-text redistribution.
